\documentclass[11pt]{article}
\usepackage[a4paper,margin=24mm]{geometry}
\usepackage{amsmath,amssymb,amsthm,mathtools}
\usepackage[hidelinks]{hyperref}
\newtheorem{theorem}{Theorem}
\newcommand{\C}{\mathbb C}
\newcommand{\R}{\mathbb R}
\newcommand{\diag}{\operatorname{diag}}
\newcommand{\ord}{\operatorname{ord}}
\title{Independent derivation of the original four-factor moment and weighted quotient formulas}
\author{Independent mathematical review}
\date{20 September 2026}
\begin{document}
\maketitle

\paragraph{Verdict and reading boundary.}
The equations OFM1--17 under review are correct on their stated admitted
five-orbit domain. In particular, no quotient correction or Gamma weight
is missing from OFM14--15. The bound $32$ follows from the four original
quadratic factors; it is a bound, not an assertion that this order is
attained by an actual period. The positive leading-moment formula is
restricted to real admitted periods, as it must be. The present proof
uses the original family and factor definitions FC5, FC10--14, NC1--5
and PCL1--10 \cite{FC,NC,PCL}; it uses the original orthogonal polynomial
recurrence and measure in WCF2--6 \cite{WCF}. It does not reprove the
entire-period construction underlying PCL or assert a conductor order
jump or an RH conclusion.

The positivity slice is already established in PCL15--16. Its application
to the original exponential curve, and the individual-factor conjugation
identity, are made explicit below. The classical orthogonal-polynomial
dictionary used by WCF credits the NIST Digital Library of Mathematical
Functions, equations (18.19.8)--(18.19.9), (18.22.8), and (18.23.7)
\cite{DLMF}. In this review the required generating function is proved
directly from the retained recurrence; no independent reading of a
different primary-source edition is claimed.

\section{Nine exact frequencies for each original factor}
Retain $0<\delta<1/2$, $\gamma>2$, and the ordered quartet
\[
 (\rho_1,\rho_2,\rho_3,\rho_4)
 =\left(\tfrac12+\delta+i\gamma,\tfrac12+\delta-i\gamma,
 \tfrac12-\delta+i\gamma,\tfrac12-\delta-i\gamma\right).
 \tag{IR1}
\]
Put $t_a(z)=e^{\rho_a z}$, $Q_0(t)=t_1t_4-t_2t_3$,
$D_5=\diag(\zeta,\zeta^2,\zeta^3,\zeta^4)$, where
$\zeta=e^{2\pi i/5}$. Let $H$ be the full original invertible period
matrix and put
\[
 M_j=H^{-1}D_5^{-j}H,\quad f_j(z)=Q_0(M_jt(z)),\quad
 E_A(z)=\prod_{j=1}^4 f_j(z),\quad
 B_0=\frac12\begin{pmatrix}0&0&0&1\\0&0&-1&0\\
 0&-1&0&0\\1&0&0&0\end{pmatrix}.
 \tag{IR2}
\]
Thus $C_j=M_j^{\mathsf T}B_0M_j$ is symmetric and
$f_j=t^{\mathsf T}C_jt$. With
$\alpha=(1,1,0,0)$ and $\beta=(1,0,1,0)$, direct addition gives
\[
 \rho_a+\rho_b
 =1+(2(\alpha_a+\alpha_b)-2)\delta
     +i(2(\beta_a+\beta_b)-2)\gamma.
\]
Consequently, writing $0\le u,w\le2$,
\[
 c_{j,uw}=\sum_{\alpha_a+\alpha_b=u,\,\beta_a+\beta_b=w}(C_j)_{ab},
 \quad \lambda_{uw}=1+(2u-2)\delta+i(2w-2)\gamma,
 \quad f_j(z)=\sum_{u,w=0}^2c_{j,uw}e^{\lambda_{uw}z}.
 \tag{IR3}
\]
This is a sum over ordered pairs $(a,b)$, so a nondiagonal quadratic
coefficient contributes twice as required. No factor $2$ has been lost.

For completeness, the kernel of restriction of homogeneous quadratics
to $t(z)$ is precisely $\C Q_0$. There are ten quadratic monomials.
Their frequency labels $(u,w)$ comprise all nine pairs in $\{0,1,2\}^2$;
the only repeated pair is $(1,1)$, arising from $t_1t_4$ and $t_2t_3$.
The nine distinct exponentials are independent: their derivatives at
zero through degree eight form the matrix
\[
 \mathcal V_{n,(u,w)}=\lambda_{uw}^n,\qquad
 \det\mathcal V=\prod_{a<b}(\lambda_b-\lambda_a)\ne0.
 \tag{IR4}
\]
The exponents are distinct because their real parts determine $u$ and
their imaginary parts determine $w$. The determinant identity follows
because this alternating polynomial of the exponents is divisible by
each difference, has the same degree as their product, and has leading
coefficient one. This also proves the exact finite-jet inverse
$c_j=\mathcal V^{-1}(f_j(0),\ldots,f_j^{(8)}(0))^{\mathsf T}$, with
$\mathcal V^{-1}=\operatorname{adj}(\mathcal V)/\det\mathcal V$.

The admitted original quadrics
$Q_0(H^{-1}X)$ and $Q_0(H^{-1}D_5^{-j}X)$ are nonassociate. If $f_j$
were identically zero, the preceding quadratic kernel statement would
give $Q_0(M_jt)=cQ_0(t)$ as polynomials. Invertibility of $M_j$ and
$B_0$ excludes $c=0$. Substitution $t=H^{-1}X$ would therefore make
those two original quadrics associates, a contradiction. Thus every
$f_j$ is nonzero and IR4 proves
\[
 n_j:=\ord_0f_j\in\{0,\ldots,8\},\qquad
 a_{j,k}:=f_j^{(k)}(0)=\sum_{u,w}c_{j,uw}\lambda_{uw}^k.
 \tag{IR5}
\]
Multiplying convergent Taylor series now gives the complete moments
and their first nonzero member:
\[
 \begin{split}
 \mu_n&=\sum_{k_1+\cdots+k_4=n}
 \frac{n!}{k_1!k_2!k_3!k_4!}\prod_{j=1}^4a_{j,k_j},\\
 v&=\sum_{j=1}^4n_j\le32,\qquad
 \mu_v=v!\prod_{j=1}^4\frac{a_{j,n_j}}{n_j!}\ne0.
 \end{split}\tag{IR6}
\]
Every summand below degree $v$ has a vanishing factor. At degree $v$
the only nonzero candidate has $k_j=n_j$ for each $j$, proving both
the order statement and the exact constant.

\section{Conjugation, the original unit, and the exact leading moment}
Write the original family as in NC2 and PCL7:
\[
 H(u)=D_uR_uV^{-1}U,\quad R_u=\mathcal R(u^{-1}),\quad
 V_{a,k}=w_a^k,\quad w_a=\rho_a-\tfrac12,\quad
 U=\diag(a,\bar a,\bar a,a),\quad a\ne0.
 \tag{IR7}
\]
The full $D_u$ includes the scalar, Gamma ray factors and branch
powers. It commutes with $D_5$. Hence, setting
$L_u=U^{-1}VR_u^{-1}$, multiplication with the actual inverse of $H$
gives
\[
 H(u)^{-1}D_5^{-j}H(u)=L_uD_5^{-j}L_u^{-1}.
 \tag{IR8}
\]
This is an equality in the retained coordinates; it makes no metric
replacement. In particular the original $D_u$ is still part of $H$
in every map for which it has not cancelled algebraically.

Let $J$ exchange indices $1,2$ and $3,4$. The explicit quartet and
unit give $\bar V=JV$, $\bar U=JUJ$, and $Q_0(Jt)=-Q_0(t)$.
Every PCL2 coefficient is real: $\beta=2(\gamma^2-\delta^2)$,
$\eta=(\delta^2+\gamma^2)^2$, all factorials, the integer powers,
and the Pochhammer factors $(b/5)_L$ are real. The convergent entire
series therefore gives $\overline{R_u}=R_{\bar u}$. As a result,
\[
 \overline{L_u}=JL_{\bar u},\qquad
 \overline{L_u^{-1}}=L_{\bar u}^{-1}J,\qquad
 \overline{t(\bar z)}=Jt(z).
 \tag{IR9}
\]
For $f^\#(z)=\overline{f(\bar z)}$, substitution of all three
identities, together with $\overline{D_5^{-j}}=D_5^j$, yields
\[
 \begin{split}
 f_{j,u}^{\#}(z)
 &=Q_0\!\left(JL_{\bar u}D_5^jL_{\bar u}^{-1}t(z)\right)\\
 &=-Q_0\!\left(L_{\bar u}D_5^{-(5-j)}L_{\bar u}^{-1}t(z)\right)
 =-f_{5-j,\bar u}(z).
 \end{split}\tag{IR10}
\]
This calculation explains the minus sign and the conjugate period.
The pairwise nonassociate domain is preserved: conjugation followed
by the invertible substitution $J$ permutes the five quadrics, with
the harmless common nonzero sign. No unproved closure of an arbitrary
period domain is needed beyond membership in that specified locus.

Put $g_u=f_{1,u}f_{2,u}$. Applying IR10 twice gives
\[
 E_{A,u}=g_u g_{\bar u}^{\#},\qquad
 \mu_n(u)=\sum_{k=0}^n\binom nk
 g_u^{(k)}(0)\overline{g_{\bar u}^{(n-k)}(0)}.
 \tag{IR11}
\]
For real admitted $u$, this is $E_{A,u}(x)=|g_u(x)|^2$ for real $x$.
If $\nu=n_1+n_2$, then IR10 gives $n_4=n_1,n_3=n_2$, and hence
\[
 v=2\nu\le32,\quad
 \mu_v=\binom{2\nu}{\nu}|g_u^{(\nu)}(0)|^2>0,\qquad
 \mu_n\in\R\quad(n\ge0).
 \tag{IR12}
\]
To check the exact factorial, the leading term of $g_u$ is
$g_u^{(\nu)}(0)z^\nu/\nu!$. Its product with its conjugate
reflection has coefficient $|g_u^{(\nu)}(0)|^2/(\nu!)^2$ at
$z^{2\nu}$, which must be multiplied by $(2\nu)!$ to obtain
the derivative. Reality of all derivatives follows either from
IR11 by replacing $k$ by $n-k$, or by conjugate reflection of
the entire function. Higher derivatives need not be positive.

On positive real periods this agrees exactly with the earlier PCL15
positivity slice. The conjugation proof above covers every admitted
real period, including any negative real period in the stated domain. Indeed
set $Z=e^{2\delta x}$, $W=e^{2i\gamma x}$ and let $A(Z,W)$ be
PCL8's original biform. Homogeneity of degree eight gives
\[
 E_A(x)=e^{(4-8\delta-8i\gamma)x}A(Z,W)
       =e^{(4-8\delta)x}W^{-4}A(Z,W)\ge0.
 \tag{IR13}
\]
Thus the source of the positivity phenomenon is retained explicitly;
IR10--12 evaluate the factor and moment statements needed here.

\section{All four quotient maps and their complete matrices}
Let $\mathcal P_m=\C[S]_{\le m}$, with $\mathcal P_{-1}=0$.
Fix the original cutoff $D\ge0$, mass $M_s=(2\pi)^{s/2}$,
$s\in\{1,a_{\rm amp}\}$, $b=s/2$, and centre $c_0=a_{\rm amp}/2$.
For $0\le j\le4$ put
\[
 r_j=\sum_{k=j+1}^4n_k,\quad c_j=c_0-j,\quad
 \mathcal H_j=\mathcal P_{D+r_j}/\mathcal P_{r_j-1}.
 \tag{IR14}
\]
The space $\mathcal H_j$ has the quotient norm induced by the full
original Gamma measure at centre $c_j$, namely
\[
 \|P\|_{s,c}^2=\int_{\R}|P(c+iy)|^2
 \frac{(2\pi)^{s/2}2^{s/2}}{4\pi\Gamma(s/2)}
 |\Gamma(s/4+iy/2)|^2\,dy.
 \tag{IR15}
\]
Define the actual factor translation by
$T_jP(S)=\sum_{u,w}c_{j,uw}P(S+\lambda_{uw})$. On monomials,
\[
 T_jS^N=\sum_{k=0}^N\binom Nk a_{j,k}S^{N-k}
       =\sum_{k=n_j}^N\binom Nk a_{j,k}S^{N-k}.
 \tag{IR16}
\]
Thus $T_j\mathcal P_{r_{j-1}-1}\subset\mathcal P_{r_j-1}$
and $T_j\mathcal P_{D+r_{j-1}}\subset\mathcal P_{D+r_j}$.
The induced quotient map is bijective, because on the ordered retained
monomials $S^{r_{j-1}+q}$ it has distinct leading degrees $r_j+q$
and nonzero diagonal coefficients
$\binom{r_{j-1}+q}{n_j}a_{j,n_j}$. Successive degree elimination
gives its inverse. This includes $n_j=0$ and $r_j=0$ without
an extra convention other than $\mathcal P_{-1}=0$.
Expansion of the four finite translation sums gives exactly
\[
 T_A=T_4T_3T_2T_1:\mathcal H_0\longrightarrow\mathcal H_4.
 \tag{IR17}
\]
The denominator of each intermediate quotient is mapped into the
next denominator. Therefore choosing a different intermediate
representative has no effect on the final image. No intermediate
low-degree component has been incorrectly discarded.

Retain the WCF2 orthogonal polynomials
\[
 p_0(y)=1,\quad p_1(y)=y,\quad
 p_{n+1}(y)=yp_n(y)-n(n-1+b)p_{n-1}(y),
 \quad h_n=M_s n!(b)_n.
 \tag{IR18}
\]
Their orthonormal versions are
$\phi_{n,c}(S)=p_n((S-c)/i)/\sqrt{h_n}$; the quotient frame
for $\mathcal H_j$ is $[\phi_{r_j,c_j}],\ldots,[\phi_{r_j+D,c_j}]$.
To derive the generating series directly, set
$F(y,t)=\sum_{n\ge0}p_n(y)t^n/n!$. Multiplying IR18 by
$t^n/n!$ and summing gives
$(1+t^2)\partial_tF=(y-bt)F$, $F(y,0)=1$.
Integration in the disk $|t|<1$, with branches having value one
at zero, gives
\[
 F(y,t)=(1+t^2)^{-b/2}e^{y\arctan t},\qquad
 \omega(t)=-i\arctan t.
 \tag{IR19}
\]
For $\psi_{n,c}(S)=p_n((S-c)/i)/n!$, translation from $c_{j-1}$
to $c_j=c_{j-1}-1$ has the multiplier
\[
 \sum_{n\ge0}(T_j\psi_{n,c_{j-1}})(S)t^n
 = e^{-\omega(t)}f_j(\omega(t))
   \sum_{m\ge0}\psi_{m,c_j}(S)t^m.
 \tag{IR20}
\]
The exponent is correct because
$(S+\lambda-c_{j-1})/i=(S-c_j)/i+(\lambda-1)/i$.
Set
\[
 G_j(t)=t^{-n_j}e^{-\omega(t)}f_j(\omega(t)),\qquad
 (G_j)_0=(-i)^{n_j}a_{j,n_j}/n_j!\ne0.
 \tag{IR21}
\]
Comparing the coefficient of $t^{r_{j-1}+q}$ in IR20, an output
frame vector of degree $r_j+p$ has coefficient
$(G_j)_{q-p}$ for $p\le q$, zero for $p>q$. Coefficients with
degree below $r_j$ represent zero in the quotient and have already
been accounted for by IR16. Finally
\[
 \phi_{n,c}=\frac{\rho_n}{\sqrt{M_s}}\psi_{n,c},\qquad
 \rho_n=\sqrt{\frac{n!}{(b)_n}},
\]
so the two unchanged factors $M_s^{-1/2}$ cancel and the full
matrix entry is
\[
 (A_j)_{pq}=(G_j)_{q-p}\frac{\rho_{r_{j-1}+q}}{\rho_{r_j+p}}
 \quad(p\le q),\qquad
 A_j=W_j^{-1}T_D(G_j)W_{j-1},\quad
 W_j=\diag(\rho_{r_j},\ldots,\rho_{r_j+D}).
 \tag{IR22}
\]
This proves the offsets directly from input and output degrees,
without inferring them from an unweighted matrix.

Finite upper triangular Toeplitz matrices multiply by truncated
convolution. The four intermediate $W_j$ cancel in order, and
\[
 \prod_{j=1}^4G_j(t)=t^{-v}e^{-4\omega(t)}E_A(\omega(t))=g(t),
 \qquad A_4A_3A_2A_1=W_4^{-1}T_D(g)W_0=A_{D,s}.
 \tag{IR23}
\]
This is precisely WCF5--6, with $r_0=v,r_4=0$, original centres
$a_{\rm amp}/2$ and $a_{\rm amp}/2-4$, and the original mass.
If $B_j(t)=\sum_{q=0}^D b_{j,q}t^q$, define
\[
 b_{j,0}=(G_j)_0^{-1},\qquad
 b_{j,q}=-(G_j)_0^{-1}\sum_{k=1}^q(G_j)_k b_{j,q-k}.
 \tag{IR24}
\]
Then every coefficient through degree $D$ of $G_jB_j-1$ is zero,
so $A_j^{-1}=W_{j-1}^{-1}T_D(B_j)W_j$ exactly. Thus
\[
 \begin{split}
 A_{D,s}^{-1}&=A_1^{-1}A_2^{-1}A_3^{-1}A_4^{-1},\\
 \|\wedge^k A_{D,s}^{-1}\|&\le
       \prod_{j=1}^4\|\wedge^k A_j^{-1}\|\quad(0\le k\le D+1),\\
 \det A_{D,s}&=
 \left(\frac{(-i)^v\mu_v}{v!}\right)^{D+1}
 \prod_{q=0}^D\frac{\rho_{v+q}}{\rho_q}.
 \end{split}\tag{IR25}
\]
Functoriality of exterior powers and the inequality
$\|BCx\|\le\|B\|\|C\|\|x\|$ prove the middle statement in
the retained induced Hilbert norms. The determinant is the product
of the triangular diagonals, with IR6 supplying the exact leading
constant. It does not bound the individual inverse exterior norms
by a uniform arithmetic constant; such a bound is not claimed here.

\section{The actual marked receiver}
In the original marking write $X,Y$ for the columns
$x_a=[S^v\ell_a(S)]$ and $y_a=T_Ax_a$, and retain the literal
invertible four-point matrix $K$ from FC15a. Then
$\Phi_*=XK^{-1}$ and $\Psi_*=YK^{-1}$. Since IR17 is an equality
of the original operators, it gives
\[
 T_4T_3T_2T_1\Phi_*q=\Psi_*q\qquad(q\in\C^4).
 \tag{IR26}
\]
This includes every affine signed-state column. It does not replace
the nonlinear signed-state transformation by a linear conductor.
If $V_\rho$ denotes FC8's matrix of powers and $C_\kappa$ is the
unchanged translation on output coefficients, direct application
of IR16 to $S^{v+q}$ yields
\[
 M_*=C_\kappa B_v(V_\rho^{\mathsf T})^{-1}K^{-1},\qquad
 (B_v)_{rq}=\binom{v+q}{v+q-r}\mu_{v+q-r}
 \quad(0\le r\le q\le3).
 \tag{IR27}
\]
All its moments are the exact finite sums IR3, IR5--6. This is the
coefficient map needed to return the moment calculation to the
marked receiver, with no independent choice of those moments.

\paragraph{Scope of the executable certificate.}
The companion Python file checks ordered-pair coefficients against
direct quadratic substitution, the nine-node determinant and inverse
jet certificate, exact conjugation in $\mathbb Q(i)[\zeta]/(\Phi_5)$,
direct translations of recurrence polynomials in all four quotient
frames at two sets of nonzero offsets and both original Gamma
parameters, their weighted product and inverse, and the factorial
in the positive leading moment. These are finite exact tests of the
identities proved above. Arbitrary real matrices and prescribed jet
test factors in that checker are explicitly auxiliary tests and
are not labelled actual finite periods. The actual coefficient limit
is identified separately. Source hashes and test counts are in the
generated JSON receipt.

\begin{thebibliography}{9}
\bibitem{FC} Split-Zero programme, \emph{The ES--Fable inverse
correspondence in the original weighted conductor},
\texttt{FABLE\_TO\_ORIGINAL\_CONDUCTOR.tex}, FC1--25 and FC26--32.
The work's human-source attributions are retained in that cumulative
source; this citation does not assign human authorship to this review.
\bibitem{NC} Split-Zero programme, \emph{Vanishing of the original
proper-source kernel through conductor coprimality},
\texttt{ORIGINAL\_CONDUCTOR\_COPRIMALITY.tex}, NC1--5.
\bibitem{PCL} Split-Zero programme, \emph{The literal original
conductor limit and an invertible pivot restriction},
\texttt{PCL\_COMPLETE.tex}, PCL1--16.
\bibitem{WCF} Split-Zero programme, \emph{A polynomial forward bound
for the full weighted conductor},
\texttt{WEIGHTED\_CONDUCTOR\_FORWARD.tex}, WCF1--6.
\bibitem{DLMF} F. W. J. Olver, A. B. Olde Daalhuis, D. W. Lozier,
B. I. Schneider, R. F. Boisvert, C. W. Clark, B. R. Miller,
B. V. Saunders, H. S. Cohl and M. A. McClain, editors,
\emph{NIST Digital Library of Mathematical Functions}, Chapter 18,
Meixner--Pollaczek weight, recurrence and generating function,
\url{https://dlmf.nist.gov/18.19},
\url{https://dlmf.nist.gov/18.22},
\url{https://dlmf.nist.gov/18.23}.
The locators and attribution here follow the already-read WCF source.
\end{thebibliography}
\end{document}
